\documentclass[conference]{IEEEtran}

\usepackage{cite}
\usepackage{amsmath,amssymb,amsfonts}
\usepackage{graphicx}
\usepackage{textcomp}
\usepackage{xcolor}
\usepackage{hyperref}
\usepackage{booktabs}
\usepackage{multirow}
\usepackage{tikz}
\usepackage{pgfplots}
\pgfplotsset{compat=1.18}

\def\BibTeX{{\rm B\kern-.05em{\sc i\kern-.025em b}\kern-.08em
    T\kern-.1667em\lower.7ex\hbox{E}\kern-.125emX}}

\begin{document}

\title{Individual Impact of Industry 4.0 and the Transition Toward Industry 5.0}

\author{\IEEEauthorblockN{1\textsuperscript{st} Thaveesak Srikuruwal}
\IEEEauthorblockA{\textit{Computer Engineering} \\
\textit{of Mae Fah Luang University}\\
Thasud, Muang, Chiang Rai \\
6871501001@lamduan.mfu.ac.th}
\and
\IEEEauthorblockN{2\textsuperscript{nd} Suttidech Jittawisuttikul}
\IEEEauthorblockA{\textit{Computer Engineering} \\
\textit{of Mae Fah Luang University}\\
Thasud, Muang, Chiang Rai \\
6871501002@lamduan.mfu.ac.th
}
}

\maketitle

\begin{abstract}
The rapid evolution of industrial technologies has significantly transformed production systems, workforce requirements, and societal structures. Industry~4.0, characterized by cyber-physical systems, automation, and data-driven intelligence, has reshaped both organizational processes and individual roles. This paper examines the historical progression from Industry~1.0 to Industry~4.0, analyzes the individual impact of Industry~4.0, and discusses the emerging vision of Industry~5.0, which emphasizes human-centric, sustainable, and resilient industrial development.
\end{abstract}

\begin{IEEEkeywords}
Industry 4.0, Industry 5.0, Industrial Revolution, Digital Transformation, Human-Centric Industry
\end{IEEEkeywords}

\section{Introduction}
The rapid advancement of digital technologies has fundamentally transformed industrial systems, leading to the emergence of the concept known as Industry 4.0. This paradigm emphasizes the integration of cyber–physical systems, the Internet of Things (IoT), cloud computing, and data analytics to enable intelligent, automated, and interconnected production environments [1], [2]. Industry 4.0 has been widely adopted across sectors such as manufacturing, agriculture, energy, and healthcare, offering improvements in efficiency, scalability, and operational transparency.

As these technologies become increasingly pervasive, critical discussions have emerged regarding their broader social, economic, and environmental implications. In response, the concept of Industry 5.0 has been proposed to complement and extend Industry 4.0 by placing stronger emphasis on human-centric design, sustainability, and system resilience [3]. This report addresses two key questions: (1) What is the individual impact of Industry 4.0? and (2) What should Industry 5.0 be? The analysis draws on existing academic literature to provide a balanced and forward-looking perspective.

\section{Evolution of Industrial Revolutions}
The industrial landscape has evolved through several distinct phases, each defined by dominant technologies and production paradigms.

Industry~1.0 began in the late 18th century with the introduction of mechanization powered by water and steam engines. Industry~2.0 followed with the adoption of electricity, enabling mass production and assembly-line manufacturing. Industry~3.0 introduced electronics, computers, and automation, significantly improving production efficiency and precision. As shown in Fig.~\ref{fig:industry_history}, Industry~4.0 builds upon these foundations by connecting physical systems with digital intelligence, enabling real-time data exchange and autonomous decision-making \cite{kagermann2013}.

\begin{figure}[h]
    \centering
    \includegraphics[width=0.7\columnwidth]{figures/industry_history_1_to_4.png}
    \caption{Evolution of Industry from 1.0 to 4.0, illustrating key technological transformations.}
    \label{fig:industry_history}
\end{figure}


\section{Individual Impact of Industry 4.0}
Industry 4.0 has significantly reshaped the role of individuals within industrial ecosystems. One major impact is the shift in skill requirements. Professionals are increasingly expected to possess interdisciplinary knowledge that combines engineering, information technology, data analysis, and cybersecurity \cite{hermann2016}. Traditional job roles are evolving toward system supervision, integration, and optimization rather than manual operation.

Continuous learning has become a defining characteristic of individual participation in Industry 4.0. Rapid technological change necessitates ongoing skill development and adaptability \cite{stock2016}. Individuals must be capable of interacting with intelligent systems, interpreting data-driven insights, and making informed decisions in complex environments. These demands encourage lifelong learning but also introduce pressure on workers to remain technologically relevant.

Industry 4.0 has also altered how individuals interact with physical and digital spaces. Remote monitoring and control enable systems to be managed across large distances, improving flexibility and reducing operational costs \cite{fao2020}. However, increased connectivity exposes individuals to cybersecurity risks, privacy concerns, and system failures, highlighting the importance of secure and resilient system design \cite{mtetwa2022}.

Industry 4.0 has a direct and tangible impact on individuals through the integration of digital technologies into everyday products and services, such as smartphones and smart retail systems. Modern mobile phones combine cloud computing, artificial intelligence, and Internet of Things connectivity, enabling individuals to perform a wide range of activities—including communication, financial transactions, navigation, remote work, and access to digital services—through a single device. This has significantly improved efficiency, convenience, and flexibility in daily life, while also introducing challenges related to privacy, cybersecurity, and work–life balance. Similarly, Industry 4.0 technologies in retail environments, exemplified by Uniqlo, employ RFID-based inventory systems, data-driven demand forecasting, and automated checkout solutions to enhance operational efficiency and customer experience. For consumers, these technologies reduce waiting times and improve product availability, while for employees they shift job roles away from repetitive manual tasks toward digital system interaction and customer-focused responsibilities. Together, these examples illustrate how Industry 4.0 reshapes individual behavior, skill requirements, and expectations by embedding intelligent, automated, and data-driven systems into everyday life.

\section{Limitations of Industry 4.0}
While Industry 4.0 has delivered substantial benefits, it also presents notable limitations. A primary concern is its strong emphasis on automation and efficiency, which may reduce meaningful human involvement and marginalize human creativity and judgment \cite{pfeiffer2017}. Excessive automation can weaken human agency and reduce opportunities for contextual decision-making.

Centralization is another critical limitation. Many Industry 4.0 implementations rely on centralized cloud infrastructures, creating single points of failure and increasing dependency on platform providers \cite{schwab2016}. In multi-stakeholder environments, centralized control may undermine transparency, accountability, and trust.

Sustainability is not always explicitly addressed in Industry 4.0 deployments. Although optimization techniques can reduce waste, large-scale digital infrastructures may increase energy consumption and environmental impact if sustainability is not integrated into system design \cite{bonilla2018}. These challenges underscore the need for an evolved industrial paradigm.

\section{Vision of Industry 5.0}
Industry 5.0 extends beyond technological efficiency by emphasizing human-centricity, sustainability, and resilience. Rather than replacing humans, Industry 5.0 promotes collaboration between humans and intelligent machines, where technology augments human expertise and creativity \cite{eu2021}. Humans remain central to decision-making, particularly in complex and ethical contexts.

Sustainability is a foundational principle of Industry 5.0. Future industrial systems should minimize environmental impact, promote efficient resource use, and support inclusive participation by small-scale stakeholders \cite{fao2020}. This perspective aligns industrial development with broader societal and environmental goals.

Resilience further distinguishes Industry 5.0 from previous paradigms. Decentralized architectures, transparent governance, and fault-tolerant system designs enhance robustness against cyber threats, infrastructure failures, and external disruptions \cite{mtetwa2022}. These characteristics are increasingly essential as industrial systems become more interconnected and critical.

\section{Reflective Perspective on Personal Impact}
From an individual perspective, Industry 4.0 has reshaped how I engage with technology, research, and problem-solving. Exposure to interconnected systems and data-driven architectures has emphasized the importance of interdisciplinary thinking and security-aware design. Rather than focusing solely on technical performance, I have become more attentive to system trust, transparency, and long-term sustainability.

The transition toward Industry 5.0 resonates with my research interests and professional values, particularly in the context of small-scale and resource-constrained environments such as smart farming. The emphasis on human-centric and resilient systems reinforces the idea that technological innovation should serve practical human needs rather than purely optimizing efficiency. This perspective has influenced my approach to designing systems that balance performance, cost, and trust.

\section{Conclusion}
Industry 4.0 has transformed industrial systems by enabling automation, connectivity, and data-driven decision-making, significantly affecting individual roles and skill requirements. While these advances offer efficiency and scalability, they also introduce challenges related to human agency, centralization, and sustainability.

Industry 5.0 addresses these limitations by prioritizing human-centric design, environmental responsibility, and system resilience. By fostering collaboration between humans and intelligent technologies, Industry 5.0 provides a complementary framework that builds upon the strengths of Industry 4.0 while mitigating its weaknesses. Together, these paradigms offer a pathway toward industrial systems that are efficient, trustworthy, and socially responsible.

\begin{thebibliography}{10}

\bibitem{kagermann2013}
H. Kagermann, W. Wahlster, and J. Helbig, \emph{Recommendations for Implementing the Strategic Initiative INDUSTRIE 4.0}, acatech, 2013.

\bibitem{monostori2014}
L. Monostori, ``Cyber-physical production systems,'' \emph{Procedia CIRP}, vol. 17, pp. 9--13, 2014.

\bibitem{eu2021}
European Commission, \emph{Industry 5.0: Towards a Sustainable, Human-Centric and Resilient European Industry}, 2021.

\bibitem{hermann2016}
M. Hermann, T. Pentek, and B. Otto, ``Design principles for Industrie 4.0 scenarios,'' in \emph{Proc. HICSS}, 2016.

\bibitem{stock2016}
T. Stock and G. Seliger, ``Opportunities of sustainable manufacturing in Industry 4.0,'' \emph{Procedia CIRP}, vol. 40, 2016.

\bibitem{fao2020}
Food and Agriculture Organization, \emph{Digital Technologies in Agriculture}, FAO, 2020.

\bibitem{mtetwa2022}
D. Mtetwa et al., ``A survey of LoRaWAN security,'' \emph{ICT Express}, vol. 8, no. 3, 2022.

\bibitem{pfeiffer2017}
S. Pfeiffer, ``The vision of Industrie 4.0 in the making,'' \emph{Work, Employment and Society}, vol. 31, no. 1, 2017.

\bibitem{schwab2016}
K. Schwab, \emph{The Fourth Industrial Revolution}, World Economic Forum, 2016.

\bibitem{bonilla2018}
J. Bonilla et al., ``Industry 4.0 and sustainability,'' \emph{Sustainability}, vol. 10, no. 12, 2018.

\end{thebibliography}

\end{document}
